% from ega1-10-fr.tex
\oldpage[I]{197}
  \item[\rm{(ii)}] La question étant locale, on peut supposer que l'on a
    une suite exacte
    $\mathscr{O}_X^m\to\mathscr{O}_X^n\to\mathscr{F}\to0$
    (\hyperref[0.5.3.2-fr]{0, 5.3.2})~; comme $\gamma^\sharp$ est
    fonctoriel, et que les foncteurs $i^*(\mathscr{F})$ et
    $\mathscr{F}_{/X'}$ sont exacts à droite (d'après (i) et
    (\hyperref[0.4.3.1-fr]{0, 4.3.1})), on a le diagramme commutatif
    \[
      \xymatrix{
        i^*(\mathscr{O}_X^m)\ar[r]\ar[d]_{\gamma^\sharp} &
        i^*(\mathscr{O}_X^n)\ar[r]\ar[d]_{\gamma^\sharp} &
        i^*(\mathscr{F})\ar[r]\ar[d]_{\gamma^\sharp} &
        0\\
        (\mathscr{O}_X^m)_{/X'}\ar[r] &
        (\mathscr{O}_X^n)_{/X'}\ar[r] &
        \mathscr{F}_{/X'}\ar[r] &
        0
      }
      \tag{10.8.8.1}\label{I.10.8.8.1-fr}
    \]
    dont les lignes sont exactes. En outre, les deux foncteurs
    $i^*(\mathscr{F})$ et $\mathscr{F}_{/X'}$ commutent aux sommes
    directes finies (\hyperref[0.3.2.6-fr]{0, 3.2.6} et
    \hyperref[0.4.3.2-fr]{4.3.2}) et on est donc ramené à démontrer notre
    assertion pour $\mathscr{F}=\mathscr{O}_X$. On a alors
    $i^*(\mathscr{O}_X)=(\mathscr{O}_X)_{/X'}=
    \mathscr{O}_{\widehat X}$
    (\hyperref[0.4.3.4-fr]{0, 4.3.4}), et $\gamma^\sharp$ est un
    homomorphisme de $\mathscr{O}_{\widehat X}$-Modules~; il suffit donc
    de vérifier que $\gamma^\sharp$ transforme la section unité de
    $\mathscr{O}_{\widehat X}$ au-dessus d'un ouvert de $X'$ en elle-même,
    ce qui est immédiat et montre donc que dans ce cas $\gamma^\sharp$ est
    l'identité.
\end{enumerate}

\begin{corollary}[10.8.9]
\label{I.10.8.9-fr}
Le morphisme d'espaces annelés $i:X_{/X'}\to X$ est plat.
\end{corollary}

Cela résulte en effet de (\hyperref[0.6.7.3-fr]{0, 6.7.3}) et de
(\hyperref[I.10.8.8-fr]{10.8.8, (i)}).

\begin{corollary}[10.8.10]
\label{I.10.8.10-fr}
Si $\mathscr{F}$ et $\mathscr{G}$ sont des $\mathscr{O}_X$-Modules
cohérents, il existe des isomorphismes canoniques fonctoriels (en
$\mathscr{F}$ et $\mathscr{G}$)
\[
  (\mathscr{F}_{/X'})\otimes_{(\mathscr{O}_X)_{/X'}}
  (\mathscr{G}_{/X'})
  \xrightarrow{\sim}
  (\mathscr{F}\otimes_{\mathscr{O}_X}\mathscr{G})_{/X'}
  \tag{10.8.10.1}\label{I.10.8.10.1-fr}
\]
\[
  (\mathscr{H}\!om_{\mathscr{O}_X}(\mathscr{F},\mathscr{G}))_{/X'}
  \xrightarrow{\sim}
  \mathscr{H}\!om_{(\mathscr{O}_X)_{/X'}}
  (\mathscr{F}_{/X'},\mathscr{G}_{/X'})
  \tag{10.8.10.2}\label{I.10.8.10.2-fr}
\]
\end{corollary}

Cela résulte de l'identification canonique de $i^*(\mathscr{F})$ et de
$\mathscr{F}_{/X'}$~; l'existence du premier isomorphisme est alors un
résultat valable pour tous les morphismes d'espaces annelés
(\hyperref[0.4.3.3.1-fr]{0, 4.3.3.1}) et celle du second est un résultat
valable pour tous les morphismes plats
(\hyperref[0.6.7.6-fr]{0, 6.7.6}), donc découle de
(\hyperref[I.10.8.9-fr]{10.8.9}).

\begin{proposition}[10.8.11]
\label{I.10.8.11-fr}
Pour tout $\mathscr{O}_X$-Module cohérent $\mathscr{F}$, le noyau de
l'homomorphisme canonique
$\Gamma(X,\mathscr{F})\to\Gamma(X',\mathscr{F}_{/X'})$ déduit de
$\mathscr{F}\to\mathscr{F}_{/X'}$ est formé des sections nulles dans un
voisinage de $X'$.
\end{proposition}

Il résulte de la définition de $\mathscr{F}_{/X'}$ que l'image canonique
d'une telle section est nulle. Réciproquement, si
$s\in\Gamma(X,\mathscr{F})$ a une image nulle dans
$\Gamma(X',\mathscr{F}_{/X'})$, il suffit de voir que tout $x\in X'$ admet
un voisinage dans $X$ dans lequel $s$ est nulle, et on peut donc se ramener
au cas où $X=\operatorname{Spec}(A)$ est affine, $A$ noethérien,
$X'=V(\mathfrak{J})$, où $\mathfrak{J}$ est un idéal de $A$, et
$\mathscr{F}=\widetilde M$, où $M$ est un $A$-module de type fini. Alors
$\Gamma(X',\mathscr{F}_{/X'})$ est le séparé complété $\widehat M$ de $M$
pour la topologie $\mathfrak{J}$-préadique, et l'homomorphisme
$\Gamma(X,\mathscr{F})\to\Gamma(X',\mathscr{F}_{/X'})$ est
l'homomorphisme canonique $M\to\widehat M$. On sait
(\hyperref[0.7.3.7-fr]{0, 7.3.7}) que le noyau de cet homomorphisme est
l'ensemble des $z\in M$ annulés par un élément de $1+\mathfrak{J}$. On a
donc $(1+f)s=0$ pour un $f\in\mathfrak{J}$~; pour tout $x\in X'$ on en
déduit $(1_x+f_x)s_x=0$, et comme $1_x+f_x$ est inversible dans
$\mathscr{O}_x$ ($\mathfrak{J}_x\mathscr{O}_x$ étant contenu dans l'idéal
maximal de $\mathscr{O}_x$), on a $s_x=0$, ce qui démontre la proposition.

\begin{corollary}[10.8.12]
\label{I.10.8.12-fr}
Le support de $\mathscr{F}_{/X'}$ est égal à
$\operatorname{Supp}(\mathscr{F})\cap X'$.
\end{corollary}

Il est clair que $\mathscr{F}_{/X'}$ est un
$(\mathscr{O}_X)_{/X'}$-Module de type fini
(\hyperref[I.10.8.8-fr]{10.8.8, (ii)}) et
(\hyperref[0.5.2.4-fr]{0, 5.2.4}),

